\section{Introduction}
LLMs achieve strong performance on a broad set of language tasks, and they can convert large volumes of text into structured summaries and judgments. Finance is a natural setting for these tools because much information arrives in text. News, regulatory filings, and earnings calls shape expectations long before they appear in prices. Recent studies show that language model outputs can capture signals from analyst reports \citep{kim2023llms} or financial news headlines \citep{lopez2023can}, guide asset allocation based on macroeconomic conditions \citep{kim2023if}, and uncover latent thematic structures \citep{lee2025theme}. These results point to a practical opportunity. A model that reads text can form qualitative views about the market and about individual firms, but portfolio management requires these views to influence concrete trading decisions. A practical allocator still faces a missing step. Portfolio optimization requires quantitative inputs, and it also requires uncertainty estimates so that unreliable signals do not dominate the allocation. LLMs can generate qualitative views about market sentiment and firm prospects, but there is no established systematic method that turns these unstructured insights into disciplined inputs for portfolio construction. Many existing uses of LLMs in finance stop at sentiment scores, rankings, or simple trading rules. These outputs can support analysis, but they do not provide an auditable route from language based judgments to the objects a portfolio optimizer must consume.

We fill this gap by using the Black-Litterman model as an interface between LLM generated views and quantitative allocation. Mean-variance optimization (MVO) is sensitive to error in expected returns and it can produce unstable weights when forecasts are noisy \citep{best1991sensitivity, merton1980estimating, chung2022effects}. The Black-Litterman framework mitigates this sensitivity by combining an equilibrium prior with explicit views, and by weighting each view by its stated uncertainty \citep{black1992global}. This is the structure needed to incorporate language based judgments into a standard optimization pipeline. We do not modify the Black Litterman update. We use it because it provides a clean mechanism to inject views into portfolio choice, provided we can supply view values and view uncertainty in a principled way.

We propose a data-driven mapping from LLM outputs to the inputs required by the Black-Litterman update. We treat an LLM as a stochastic return forecaster. We query the model repeatedly for the same asset and horizon under a fixed prompt and a strictly time stamped information set. This produces a distribution of numerical forecasts. We set the view to the average forecast, and we set view uncertainty to the dispersion of the forecasts. The Black-Litterman update then assigns less influence to views with high dispersion and more influence to views with low dispersion. We interpret dispersion as model uncertainty about the view rather than sampling error of the average forecast, so uncertainty does not shrink with the number of queries.

This approach operationalizes LLM capabilities for portfolio optimization without adding new modeling assumptions inside the allocator. By deriving both return forecasts and confidence levels directly from observable model outputs, the proposed method ensures an auditable route from text to portfolio weights. The approach is inherently scalable because the same procedure applies across different assets and LLMs. It also supports evaluation that matches financial practice because strict information sets prevent look-ahead contamination. In our out-of-sample backtest on a large-cap equity universe, portfolios that use this mapping within the Black-Litterman framework improve risk-adjusted performance relative to equal weight and standard mean-variance baselines after accounting for transaction costs.

%Mean–variance optimization (MVO) provides a canonical normative framework for portfolio choice, but it is fragile in practice. Small errors in expected returns generate extreme and unstable weights, undermining implementability \citep{best1991sensitivity, merton1980estimating, chung2022effects}. The Black–Litterman (BLM) model addresses this instability by combining a market‑equilibrium prior with investor views, recentering mean estimates and regularizing portfolios \citep{black1992global}. In BLM, the influence of a view depends on its level $\mathbf{q}$ and its confidence matrix $\mathbf{\Omega}$. Views with lower uncertainty receive more weight. Accordingly, the key practical question is how to specify and calibrate views and their confidence. Specifying and calibrating views has been the practical bottleneck. Standard implementations set $\mathbf{q}$ by subjective judgment and map confidence to $\mathbf{\Omega}$ using ad hoc rules \citep{he2002intuition, idzorek2007step}, choices that are hard to audit, difficult to scale, and susceptible to behavioral biases. Meanwhile, large language models (LLMs) extract forward‑looking signals from unstructured text—news, filings, and earnings calls—\citep{lopez2023can, kim2023if}. However, applications often stop at sentiment scoring or simple long–short sorts, without entering portfolio construction. The literature lacks an econometrically grounded mapping from noisy LLM predictions to the $\mathbf{q}$ and $\mathbf{\Omega}$ required by BLM. We bridge this gap with a parsimonious mapping that allows the forecasting model to quantify its own uncertainty. For each asset and forecast horizon, we query an LLM $N$ times with a fixed prompt and record $N$ stochastic draws. We set $q_i$ to the mean of the draws and set confidence to their dispersion—$\Omega_{ii}$ equals the variance for absolute views, and $\mathbf{\Omega}$ equals the full covariance matrix for joint views. The LLM’s predictive dispersion determines view confidence: high variance is automatically down‑weighted in the BLM update, whereas low variance receives greater weight. Importantly, we treat $\mathbf{\Omega}$ as model uncertainty rather than sampling error of the mean. Therefore we do not divide by $N$. This mean–variance mapping replaces subjective confidence inputs with a transparent, data‑driven procedure that scales across assets, models, and time.